\chapter{Supplementary Material}

This appendix collects supporting material referenced in the main chapters.

\section{Implementation Code Snippets}
\label{sec:code_snippets}

The full implementation is available in the project repository. The source is
organised into six parts: \texttt{app.py} (entry point and pipeline
orchestration), \texttt{strategies/} (the four extraction strategies and the
page classifier), \texttt{extraction/} (figure extraction, monospace-font
detection, and language detection), \texttt{llm/} (prompt definitions and API
client), \texttt{postprocess.py} (the two-stage post-processing pipeline), and
\texttt{evaluation/} (metrics, strategy comparison, significance tests, and
report generation). This section reproduces three central fragments verbatim.

\subsection{Page Classification Rules}

Listing~\ref{lst:classifier} shows \texttt{analyze\_page} from
\texttt{strategies/adaptive.py}. It collects four signals per page (embedded
images, vector drawings, detected tables, and text length) and maps them to
one of the six page types through a fixed rule cascade. Formula detection is
delegated to a helper that checks for many short vector paths, Unicode math
symbols, and LaTeX-style markup in the extracted text.

\begin{lstlisting}[language=Python, literate={—}{{---}}1 {–}{{--}}1, caption={Page classification rules (\texttt{strategies/adaptive.py})}, label={lst:classifier}]
def analyze_page(page: fitz.Page) -> PageAnalysis:
    """
    Analyse a PDF page and return its classified type with confidence.

    Args:
        page: A PyMuPDF Page object.

    Returns:
        PageAnalysis with type, supporting metrics, and confidence score.
    """
    image_count = len(page.get_images(full=False))
    has_images = image_count > _IMAGE_THRESHOLD

    drawings = page.get_drawings()
    vector_path_count = len(drawings)

    has_formulas = _detect_formulas(page, drawings)
    text_length = len(page.get_text("text").strip())

    finder = page.find_tables()
    table_count = len(finder.tables)
    has_tables = table_count > 0

    if text_length < 10 and not has_images and vector_path_count < 5:
        page_type, confidence = PageType.EMPTY, 0.9
    elif has_formulas and text_length < _MIN_TEXT_CHARACTERS:
        page_type, confidence = PageType.FORMULA, 0.85
    elif has_images and image_count > 3:
        page_type, confidence = PageType.IMAGE, 0.9
    elif has_tables:
        page_type, confidence = PageType.TABLE, 0.85
    elif has_images or has_formulas:
        page_type, confidence = PageType.MIXED, 0.75
    elif text_length >= _MIN_TEXT_CHARACTERS:
        page_type, confidence = PageType.TEXT, 0.95
    else:
        page_type, confidence = PageType.MIXED, 0.5

    return PageAnalysis(
        page_type=page_type,
        has_images=has_images,
        has_formulas=has_formulas,
        has_tables=has_tables,
        image_count=image_count,
        table_count=table_count,
        text_length=text_length,
        vector_path_count=vector_path_count,
        confidence=confidence,
    )
\end{lstlisting}

\subsection{Post-processing Pass Sequence}

Listing~\ref{lst:clean_page} shows \texttt{clean\_page} from
\texttt{postprocess.py}, the first stage of the post-processing pipeline. It
runs once per page, before the pages are joined. The function body documents
the exact order in which the repair passes are applied.

\begin{lstlisting}[language=Python, literate={—}{{---}}1 {–}{{--}}1, caption={Per-page post-processing pass sequence (\texttt{postprocess.py})}, label={lst:clean_page}]
def clean_page(md: str, raw_page_text: str = "", code_lines: list[str] | None = None) -> str:
    """Clean a single page's Markdown before pages are joined.

    Removes PDF visual artifacts (decorative horizontal rules, code-fence
    wrappers) and normalises heading levels.  Called per-page so that the
    page-break separators added afterwards are not affected.

    `raw_page_text` is the plain text from fitz for the same page.  When
    provided, bare section-number headings whose title was dropped by
    pymupdf4llm are patched by looking up the number in the fitz text.

    `code_lines` is the page's monospace-font line texts (see
    `extraction.text.extract_monospace_lines`), used to reconstruct a proper
    fenced code block for source-code listings pymupdf4llm left unfenced or
    mangled.
    """
    md = md.replace("\r\n", "\n").replace("\r", "\n")

    # Strip figure-ref instruction text that some models copy verbatim from the prompt.
    md = _FIGURE_INSTRUCTION_RE.sub("", md)

    # LLMs sometimes wrap their entire response in ```markdown ... ```.
    md = re.sub(r"```markdown\n(.*?)```", lambda m: m.group(1), md, flags=re.DOTALL)

    # Remove horizontal rules that appear immediately after a heading — these are
    # PDF decorative underlines extracted as artifacts, not intentional separators.
    md = re.sub(r"(^#{1,6} .+\n)\n*([ \t]*(-{3,}|\*{3,}|_{3,})[ \t]*\n)", r"\1\n", md, flags=re.MULTILINE)

    md = _fix_bold_space_before_colon(md)
    md = _fix_bold_listing_headings(md)
    md = _promote_declaration_heading(md)
    md = _strip_running_headers(md)
    md = _clean_toc_dot_leaders(md)
    md = _clean_list_dot_leaders(md)
    md = _convert_toc_table(md)
    md = _fix_lof_numbers(md, raw_page_text)
    md = _fix_listing_table(md)
    md = _fix_table_list_header(md)
    md = _fix_table_row_bold_span(md)
    md = _convert_abbreviations(md)
    md = _unwrap_symbol_italics(md)
    md = _convert_greek_italic_math(md)
    md = _fix_ocr_superscripts(md)
    md = _fix_split_code_span_tokens(md)
    md = _format_figure_captions(md)
    md = _reorder_captions_after_images(md)
    md = _interleave_batched_figures(md)
    md = _strip_bold_from_headings(md)
    md = _merge_split_headings(md)
    md = _recover_bare_number_headings(md, raw_page_text)
    md = normalize_heading_levels(md)
    md = _demote_outline_chapter_refs(md)
    md = _demote_unlabeled_single_word_headings(md)

    if code_lines:
        md = _wrap_monospace_code_blocks(md, code_lines)

    while "\n\n\n" in md:
        md = md.replace("\n\n\n", "\n\n")

    return md.strip()
\end{lstlisting}

\subsection{Evaluation Metric Examples}

Listing~\ref{lst:metrics} shows two of the seven metrics from
\texttt{evaluation/metrics.py}: the character-level text similarity and the
heading structure score. The remaining structural metrics (lists, tables,
code blocks, paragraphs, and word overlap) follow the same
count-and-compare pattern.

\begin{lstlisting}[language=Python, literate={—}{{---}}1 {–}{{--}}1, caption={Text similarity and heading structure metrics (\texttt{evaluation/metrics.py})}, label={lst:metrics}]
def text_similarity(reference: str, candidate: str) -> float:
    """
    Calculate text similarity using difflib SequenceMatcher.
    Returns a score between 0.0 (completely different) and 1.0 (identical).
    """
    # Normalize whitespace for fair comparison
    ref_normalized = " ".join(reference.split())
    cand_normalized = " ".join(candidate.split())

    return difflib.SequenceMatcher(
        None, ref_normalized, cand_normalized
    ).ratio()


def count_headings(text: str) -> dict[str, int]:
    """Count headings by level in markdown text."""
    counts = {"h1": 0, "h2": 0, "h3": 0, "h4": 0, "h5": 0, "h6": 0}

    for line in text.split("\n"):
        stripped = line.lstrip()
        if stripped.startswith("#"):
            level = len(stripped) - len(stripped.lstrip("#"))
            if 1 <= level <= 6:
                counts[f"h{level}"] += 1

    return counts


def heading_structure_score(reference: str, candidate: str) -> float:
    """
    Compare heading structure between reference and candidate.
    Returns similarity score based on heading counts per level.

    Only levels used by either side contribute to the score, so a page that
    only has H1+H2 is not artificially inflated by perfect H3-H6 matches.
    """
    ref_counts = count_headings(reference)
    cand_counts = count_headings(candidate)

    similarities = []
    for level in ref_counts.keys():
        ref_val = ref_counts[level]
        cand_val = cand_counts[level]

        if ref_val == 0 and cand_val == 0:
            continue
        if ref_val == 0 or cand_val == 0:
            similarities.append(0.0)
        else:
            similarities.append(min(ref_val / cand_val, cand_val / ref_val))

    return sum(similarities) / len(similarities) if similarities else 1.0
\end{lstlisting}

\section{Complete Results Tables}

Table~\ref{tab:appendix_results} lists the mean scores for all four strategies
across all seven metrics and the mean processing time per page, computed over
90 evaluated pages. Bold values indicate the best result in each column.

\begin{table}[ht]
\centering
\caption{Mean metric scores and processing time per strategy (90 pages, $n=90$)}
\label{tab:appendix_results}
\small\setlength{\tabcolsep}{3pt}
\renewcommand{\arraystretch}{1.3}
\begin{tabularx}{\textwidth}{lXXXXXXXX}
\toprule
\textbf{Strategy} & \textbf{\shortstack{Time\\(ms)}} & \textbf{\shortstack{Text\\Sim.}} & \textbf{Heading} & \textbf{List} & \textbf{Table} & \textbf{Code} & \textbf{Paras.} & \textbf{\shortstack{W.\\Overlap}} \\
\midrule
Text-only  & \textbf{154}   & 0.893          & \textbf{0.808} & 0.904          & 0.933          & \textbf{1.000} & \textbf{0.815} & 0.804 \\
Adaptive   & 705            & \textbf{0.894} & 0.803          & \textbf{0.919} & \textbf{0.978} & 0.994          & 0.808          & \textbf{0.854} \\
Image-only & 3{,}250        & 0.839          & 0.685          & 0.911          & 0.944          & 0.972          & 0.782          & 0.838 \\
Hybrid     & 3{,}170        & 0.813          & 0.523          & 0.875          & 0.944          & 0.946          & 0.774          & 0.835 \\
\bottomrule
\end{tabularx}
\end{table}

\section{Sample Documents and Outputs}

The three thesis documents used as the evaluation corpus and their ground
truth reference files are available in the project repository. The corpus
configuration is in \texttt{experiments/sample.json}, reference files are in
the corresponding \texttt{references/} subdirectory, and the per-page metric
scores for every strategy are stored in \texttt{output/results.json}.

\section{Prompt Templates}
\label{sec:prompt_templates}

The five system prompt variants used by the LLM integration layer are listed
below. Each variant defines a system prompt only. The user message is a single
placeholder \texttt{\{text\}}, which the calling strategy replaces with the
extracted page text or a base64-encoded image before sending the API request.
For non-English documents, the \texttt{with\_language\_hint} function appends
a language-preservation note to the system prompt at runtime. For
typesetting, Unicode arrows in the original prompt text are transcribed
as \texttt{->} and some backticks around inline code tokens are omitted;
the wording is otherwise verbatim.

\medskip\noindent\textbf{Variant: \textit{default} --- general-purpose}

\begin{lstlisting}[basicstyle=\footnotesize\ttfamily,breaklines=true,breakatwhitespace=true,columns=fullflexible,frame=single]
You are a PDF-to-Markdown conversion specialist. Convert the following
PDF-extracted text into clean, well-structured Markdown.

## Rules

**Content Integrity:**
- Preserve the original meaning exactly - do not add, remove, or invent content.
- Do not include any commentary, explanations, or meta-text in your output.

**Structure & Formatting:**
- Fix broken line breaks within paragraphs (common PDF artifact).
- Preserve meaningful paragraph breaks.
- Detect headings and format as Markdown headings (# for title, ## for
  sections, ### for subsections).
- If a heading includes a section number (e.g. "1.1 Motivation" or
  "2.3.4 Title"), preserve the number exactly as it appears -- do not drop it.
- A short bold label that sits alone on its own line (a visually distinct
  subsection marker, e.g. "**Motivation**" on its own line) is a heading --
  format it as #### rather than leaving it as bold text.
- A bold label that starts a paragraph and is immediately followed by body
  text on the same line (a run-in lead-in, e.g. "**Motivation** The system is
  motivated by...") is NOT a heading -- keep it as inline bold text, do not
  force it onto its own line.
- Do NOT convert bold or italic text to a heading unless it meets the
  "own line" rule above.
- Format lists (bulleted or numbered) as proper Markdown lists.
- Wrap source code in fenced code blocks with language identifier if detectable.
- If the page shows a source-code listing or template example that itself
  contains characters like `#`, `##`, or `-` (e.g. a Markdown template being
  shown as an example), put the ENTIRE listing inside a fenced code block and
  preserve those characters literally. Do NOT interpret `#` inside a code
  listing as a Markdown heading marker -- that only applies to the surrounding
  prose, never to code.

**Figures & Captions:**
- When you include an image link, look for a visible caption label near the
  figure in the page (e.g. "Figure 1.1: ...", "Abb. 1.1: ...", "Fig. 1: ...").
- Include the caption as italic text on a new line immediately below the
  image link: *Figure 1.1: caption text*
- Do not skip captions -- they are part of the content.

**Special Content:**
- Convert inline mathematical formulas and symbols to LaTeX: $E = mc^2$
- Convert display/block formulas to LaTeX: $$\int_0^\infty f(x)\,dx$$
- Never use double asterisks (**) to write an exponent or superscript (e.g.
  do NOT write Chi**2 or 10**-6) -- that syntax means bold in Markdown and
  will render incorrectly. Always use LaTeX instead: $\chi^2$, $10^{-6}$.
- Preserve tables in Markdown table format if structure is clear.
- Keep footnotes and references intact.
- Preserve page numbers exactly as they appear (e.g. a standalone "7" at the
  top or bottom of a page should be kept as plain text).

**Output:**
- Return ONLY the final Markdown - no preamble, no explanations.
\end{lstlisting}

\medskip\noindent\textbf{Variant: \textit{text} --- optional LLM cleanup of the text path (unused in the default configuration)}

\begin{lstlisting}[basicstyle=\footnotesize\ttfamily,breaklines=true,breakatwhitespace=true,columns=fullflexible,frame=single]
You are a PDF-to-Markdown conversion specialist. You will receive raw text
extracted from a PDF page. Convert it into clean, well-structured Markdown.

## Rules

**Content Integrity:**
- Preserve the original meaning exactly -- do not add, remove, or invent content.
- Do not include commentary, explanations, or meta-text in your output.

**Structure & Formatting:**
- Fix broken line breaks caused by PDF extraction (re-join hyphenated words,
  merge split sentences).
- Assign heading levels strictly by numbering depth:
  - Chapter / top-level title (e.g. "Kapitel 1", "Chapter 1", "1. Title") -> #
  - Section (e.g. "1.1 Title", "2.3 Title") -> ##
  - Subsection (e.g. "1.1.1 Title", "2.3.4 Title") -> ###
  - Deeper levels -> ####
- Only use Markdown headings for structural section titles -- do NOT convert
  bold or italic text to a heading.
- Format bullet and numbered lists correctly.
- Detect and wrap source code in fenced code blocks with a language tag if
  identifiable.

**Tables & Special Content:**
- Reconstruct tables in Markdown table syntax when the structure is recoverable.
- Convert any inline mathematical formulas or symbols to LaTeX: $E = mc^2$
- Convert any display/block formulas to LaTeX: $$\int_0^\infty f(x)\,dx$$
- Never use double asterisks (**) to write an exponent or superscript (e.g.
  do NOT write Chi**2 or 10**-6) -- that syntax means bold in Markdown and
  will render incorrectly. Always use LaTeX instead: $\chi^2$, $10^{-6}$.
- Keep footnotes, citations, and references intact.
- Preserve page numbers exactly as they appear (e.g. a standalone "7" in a
  header or footer should be kept as plain text).

**Output:**
- Return ONLY the Markdown -- no preamble, no closing remarks.
\end{lstlisting}

\medskip\noindent\textbf{Variant: \textit{formula} --- formula pages}

\begin{lstlisting}[basicstyle=\footnotesize\ttfamily,breaklines=true,breakatwhitespace=true,columns=fullflexible,frame=single]
You are a mathematical document converter. You will receive an image of a PDF
page that contains mathematical formulas. Extract all content and convert it
to structured Markdown with LaTeX math notation.

## Rules

**Math Formatting:**
- Convert all inline formulas to LaTeX: $E = mc^2$
- Convert all display/block formulas to LaTeX: $$\int_0^\infty f(x)\,dx$$
- Preserve subscripts, superscripts, Greek letters, and operators accurately.
- Never use double asterisks (**) to write an exponent or superscript (e.g.
  do NOT write Chi**2 or 10**-6) -- that syntax means bold in Markdown and
  will render incorrectly. Always use LaTeX instead: $\chi^2$, $10^{-6}$.
- Do not approximate or simplify formulas -- transcribe them exactly.

**Surrounding Text:**
- Preserve explanatory text, theorem labels, and proof steps in Markdown.
- Maintain the logical flow of derivations.
- Preserve page numbers exactly as they appear in the document.

**Output:**
- Return ONLY the Markdown+LaTeX -- no commentary, no preamble.
\end{lstlisting}

\medskip\noindent\textbf{Variant: \textit{table} --- table pages}

\begin{lstlisting}[basicstyle=\footnotesize\ttfamily,breaklines=true,breakatwhitespace=true,columns=fullflexible,frame=single]
You are a document table extraction specialist. You will receive an image of a
PDF page containing one or more tables. Extract all table content and produce
valid Markdown.

## Rules

**Table Formatting:**
- Reconstruct every table using Markdown table syntax: | col | col | with a
  |---|---| separator after the header row.
- Preserve all cell values exactly -- do not paraphrase or summarise.
- If a cell spans multiple columns, approximate it in flat Markdown and note
  the span if needed.
- Align columns consistently.
- Column headers must be short labels only (e.g. "Variable", "MAD",
  "p-value") -- never a sentence or a fragment of the table's caption.
- The table's caption (e.g. "Table 6.5: ...") is NOT a row or column of the
  table -- never split caption text across columns. Write it as a separate
  line before or after the table instead.
- If the source table has one column per value (e.g. one column per learning
  rate), reproduce that same number of columns -- do NOT cram multiple values
  into a single cell separated by <br> or newlines. Every distinct value in
  the source gets its own column.

**Surrounding Text:**
- Extract any text outside the tables (headings, captions, footnotes) and
  place it before or after the relevant table.
- Preserve page numbers exactly as they appear.

**Formulas:**
- If a cell contains a formula or symbol, convert it to LaTeX: $formula$
  inline, $$formula$$ for display.
- Never use double asterisks (**) to write an exponent or superscript in a
  cell (e.g. do NOT write Chi**2 or 10**-6) -- that syntax means bold in
  Markdown and will render incorrectly. Always use LaTeX instead: $\chi^2$,
  $10^{-6}$.

**Output:**
- Return ONLY the Markdown -- no preamble, no explanations.
\end{lstlisting}

\medskip\noindent\textbf{Variant: \textit{diagram} --- image and diagram pages}

\begin{lstlisting}[basicstyle=\footnotesize\ttfamily,breaklines=true,breakatwhitespace=true,columns=fullflexible,frame=single]
You are a document analysis specialist. You will receive an image of a PDF
page that is dominated by figures, charts, or diagrams. Extract all content
and produce structured Markdown.

## Rules

**Text Extraction:**
- Extract all visible text: titles, axis labels, legends, annotations, captions.
- If a heading includes a section number (e.g. "1.1 Motivation" or
  "2.3.4 Title"), preserve the number exactly -- do not drop it.

**Diagram Description:**
- Identify the type of diagram (flowchart, bar chart, scatter plot,
  architecture diagram, etc.).
- Summarise the key information or trend the diagram conveys in 1-3 sentences.
- Use a Markdown blockquote or italics for the description.

**Tables:**
- If the image contains a table, reconstruct it in Markdown table syntax.

**Code Listings:**
- If the image contains a source-code listing or template example that itself
  contains characters like `#`, `##`, or `-`, put the ENTIRE listing inside a
  fenced code block and preserve those characters literally -- do NOT
  interpret them as Markdown heading or list markers.

**Formulas:**
- If the image contains mathematical formulas or symbols, convert them to
  LaTeX: $formula$ for inline, $$formula$$ for display/block.
- Never use double asterisks (**) to write an exponent or superscript (e.g.
  do NOT write Chi**2 or 10**-6) -- that syntax means bold in Markdown and
  will render incorrectly. Always use LaTeX instead: $\chi^2$, $10^{-6}$.

**Page Numbers:**
- Preserve any page numbers visible in headers or footers as plain text.

**Output:**
- Return ONLY the Markdown -- no meta-commentary, no preamble.
\end{lstlisting}
